\chapter{数值解法简介}
\intro{本章介绍有限差分法和有限元方法的基本思想。}

\section{有限差分法}

\begin{mydefinition}{差分格式}
用差商近似微商：
\begin{align*}
\frac{\partial u}{\partial t}&\approx\frac{u^{n+1}_i-u^n_i}{\Delta t}\\
\frac{\partial^2 u}{\partial x^2}&\approx\frac{u^n_{i+1}-2u^n_i+u^n_{i-1}}{(\Delta x)^2}
\end{align*}
\end{mydefinition}

\begin{theorem}[稳定性条件]
热传导方程显式格式的稳定性条件（CFL 条件）：
\[
r=\frac{k\Delta t}{(\Delta x)^2}\leq\frac{1}{2}
\]
\end{theorem}

\section{有限元方法}

\begin{mydefinition}{有限元方法}
将区域剖分为有限个单元，在每个单元上用简单函数（如线性函数）近似真实解。
\end{mydefinition}

\begin{remark}
有限差分法实现简单，适用于规则区域；有限元方法适用于复杂几何，是工程仿真（如 ANSYS）的核心算法。
\end{remark}
